\chapter{Metodología y especificación de requisitos}
\label{chap:metodologia}

\section{Enfoque metodológico}

El desarrollo siguió un enfoque guiado por especificaciones. La propuesta delimitó el problema; las especificaciones expresaron comportamiento observable; el diseño documentó decisiones y compromisos; las tareas dividieron la construcción en unidades verificables. Los términos \emph{debe}, \emph{deberá} y \emph{puede} se interpretan como niveles normativos siguiendo la convención de \textcite{rfc2119}.

El proceso se organizó en cinco fases:

\begin{enumerate}
  \item \textbf{Especificación:} convertir necesidades en requisitos identificables y escenarios Dado--Cuando--Entonces.
  \item \textbf{Diseño:} relacionar cada requisito con componentes, modelos de datos y decisiones de consistencia.
  \item \textbf{Implementación incremental:} se estableció primero la infraestructura de pruebas y se construyeron cortes verticales pequeños.
  \item \textbf{Verificación:} se ejecutaron pruebas automatizadas y experimentos reproducibles contra criterios de aceptación.
  \item \textbf{Documentación de resultados:} se definieron comandos, salidas estables, capturas pertinentes y análisis de desviaciones.
\end{enumerate}

OpenSpec constituye la línea base del proyecto a la fecha de este informe \parencite{proyectoOpenSpec}. Si la implementación exige modificar un requisito, el cambio deberá actualizar primero esa línea base y luego la matriz de trazabilidad del apéndice~\ref{app:trazabilidad}.

\section{Actores}

\begin{table}[H]
  \centering
  \caption{Actores y responsabilidades}
  \label{tab:actores}
  \begin{tabularx}{\textwidth}{@{}p{3cm}YY@{}}
    \toprule
    \textbf{Actor} & \textbf{Objetivo} & \textbf{Permisos} \\
    \midrule
    Invitado & Consultar el catálogo y registrar un pedido & Lectura de catálogo activo y creación de pedidos \\
    Operador & Administrar la operación del restaurante & Inicio de sesión, catálogo, estados, proyecciones y reproducción \\
    Trabajador & Materializar eventos pendientes & Reclamo de salida, escritura idempotente y actualización de estado \\
    Evaluador & Comprobar las propiedades técnicas & Ejecución local de semillas, pruebas y consultas de evidencia \\
    \bottomrule
  \end{tabularx}
\end{table}

\section{Requisitos funcionales}

\begin{longtable}{@{}p{1.4cm}p{4cm}p{8.2cm}@{}}
  \caption{Requisitos funcionales}\label{tab:rf}\\
  \toprule
  \textbf{ID} & \textbf{Nombre} & \textbf{Criterio verificable} \\
  \midrule
  \endfirsthead
  \toprule
  \textbf{ID} & \textbf{Nombre} & \textbf{Criterio verificable} \\
  \midrule
  \endhead
  RF-01 & Catálogo público & El invitado consulta únicamente artículos activos con sus precios vigentes. \\
  RF-02 & Compra calculada & La API obtiene precios del catálogo, valida cantidades, persiste instantáneas y descarta totales enviados por el cliente. \\
  RF-03 & Idempotencia de compra & La repetición de una clave con el mismo contenido devuelve el pedido original; con contenido distinto responde conflicto. \\
  RF-04 & Flujo de estados & Solo se permiten las transiciones definidas en la figura~\ref{fig:estados}. \\
  RF-05 & Acceso del operador & Las mutaciones administrativas y las consultas operacionales requieren una sesión válida. \\
  RF-06 & Proyección de eventos & Cada evento de salida se aplica a las dos vistas Cassandra y después se marca como procesado. \\
  RF-07 & Consultas particionadas & La línea temporal usa el pedido y la actividad diaria usa restaurante y fecha como claves de partición. \\
  RF-08 & Recuperación & El operador consulta retraso, errores e intentos, y puede solicitar la reproducción de eventos seleccionados. \\
  \bottomrule
\end{longtable}

\section{Requisitos no funcionales}

\begin{longtable}{@{}p{1.4cm}p{3.6cm}p{8.6cm}@{}}
  \caption{Requisitos no funcionales}\label{tab:rnf}\\
  \toprule
  \textbf{ID} & \textbf{Atributo} & \textbf{Criterio verificable} \\
  \midrule
  \endfirsthead
  \toprule
  \textbf{ID} & \textbf{Atributo} & \textbf{Criterio verificable} \\
  \midrule
  \endhead
  RNF-01 & Integridad & Los validadores e índices rechazan tipos inválidos y duplicados aun ante escritura directa. \\
  RNF-02 & Atomicidad & Una falla forzada en la compra no deja ni pedido ni evento de salida. \\
  RNF-03 & Idempotencia & Los reintentos no generan más de un pedido ni más de un evento lógico por vista. \\
  RNF-04 & Observabilidad & El estado de proyección expone cantidades, retraso más antiguo, intentos y último error. \\
  RNF-05 & Reproducibilidad & Una semilla fija produce agregaciones y evidencia estructural equivalentes en ejecuciones sucesivas. \\
  RNF-06 & Seguridad & Las credenciales no se almacenan en el código; la cookie JWT dura 15 minutos y usa \code{HttpOnly}, \code{Secure}, \code{SameSite=Strict} y \code{Path=/}. \\
  RNF-07 & Recursos & Los cuatro servicios están configurados con límites de 96, 384, 1536 y 3072~MiB, para un total de 5088~MiB. \\
  RNF-08 & Persistencia local & Los datos de ambos motores sobreviven al reinicio de sus contenedores mediante volúmenes con nombre. \\
  \bottomrule
\end{longtable}

\section{Criterios globales de éxito}

La evaluación se considera satisfactoria cuando: (a) el entorno inicia de forma determinista con cuatro servicios; (b) la compra confirma atómicamente un pedido y su evento; (c) la evidencia identifica validadores, índices y planes requeridos; y (d) las dos vistas Cassandra convergen después de duplicados e interrupciones. Estos criterios no sustituyen las pruebas detalladas del capítulo~\ref{chap:validacion}; funcionan como condiciones de salida del proyecto.
